\scanpage{42}
% archive foliation: 41

Si on se donne la \deleted{gerbe} \inserted{catégorie de Deuring \unsure{rigidifiée|associée}}
$(C'(A'),T')$ sur $A'=\prod_{\ell\ne p,\infty}\mathbb Z_\ell$ comme associée à la gerbe des racines carrées de $A'(1)$ (définie par le choix d'une clôture algébrique $\overline{\mathbb F}_p$ de $\mathbb F_p$), et la catégorie de Deuring $C(\mathbb Z_p)$ comme celle des $F$-cristaux sur $\overline{\mathbb F}_p$ (à iso) qui sont de rang $2$, \deleted{\unsure{tous}} pente $\frac12$, \inserted{avec le foncteur de Tate}, et la \deleted{gerbe}
catégorie $C(\mathbb Z_\infty)=C(\mathbb R)$ comme associée à la gerbe type $\unsure{\mathfrak g}_\infty$ de groupe $\mu_2$ sur $\mathbb R$, on trouve un \unsure{torseur}
\deleted{\unsure{\mathcal T}(C)} un certain torseur \inserted{$\mathcal T'(C)$} sous le groupe
\[
\deleted{\unsure{\mathbb Z_p}}\quad
A'^*/A'^{*2}\times
(\mathbb Z_p^*/\mathbb Z_p^{*2}\times\mathbb Z/2\mathbb Z)
\times\mathbb R^*/\mathbb R^{*2}=\Gamma'.
\]

Le groupe \inserted{$\mathbb Z/2\mathbb Z$} des classes d'isomorphie d'autoéquivalences de \deleted{$C$} opère sur ce torseur, de façon compatible avec son opération triviale sur lui-même, donc par translations, et il opère donc via un hom. \inserted{$f$ (\unsure{l'hyp. scindée})} de $\mathbb Z/2\mathbb Z$ dans ce groupe,
\deleted{\unsure{dont les}} l'image du \inserted{gén.} $f$ de $\mathbb Z/2\mathbb Z$ dans les $3$ facteurs étant respectivement
\[
(p,\ \unsure{\beta}=(0,1),\ p)
\]
\margin{N.B. \inserted{Comme $p\in\mathbb R^{*2}$}
\deleted{\unsure{on peut donc}, \unsure{elle est négligeable}, \unsure{dans}}
l'ens. des \unsure{plexions},
\deleted{\unsure{distinguer canoniquement le}}
$\unsure{\mathrm{Pl}}(C)$ il y a un système transitif d'isomorphismes!}

Considérons le quotient
\[
\Gamma=\Gamma'/f(\pm1),
\]
\deleted{On \unsure{voit donc}}
et soit $\mathcal T(C)$ le quotient de $\mathcal T'(C)$ par $\pm1$, qui est donc un torseur sous $\Gamma$.
Ce torseur, à isom unique près, ne dépend pas de $C$, \inserted{on peut l'appeler $\mathcal T$ ou $\mathcal T(\overline{\mathbb F}_p)$,}
D'ailleurs si $\Gamma'_0$ est le ss-grp de $\Gamma'$
\[
A'^*/A'^{*2}\times\mathbb Z_p^*/\mathbb Z_p^{*2}
\times\mathbb R^*/\mathbb R^{*2}
=H^1_{\deleted{\unsure{\mathrm{ét}}}}(\mathbb A,\mu_2)
\]
\margin{adèles}
on trouve que $\Gamma'\simeq\Gamma'_0\times f(\pm1)$, donc $\Gamma\simeq\Gamma'_0$ (can., donc
